\section{Limitations}\label{sec:limitations}
CRR is intended as a reliability framework, not as a universal accuracy replacement. First, the method is not an MVTec AUROC SOTA claim. Official SubspaceAD remains a very strong ranking baseline, and our MVTec contribution is mainly storage, calibration, and diagnostics. Second, CRR is not adversarially robust: label-aware PCA-residual FGSM tests show severe collapse and objective sensitivity. Third, conformal reliability depends on whether the few normal support images are representative of future normal data. If support normals are not exchangeable with test normals, p-values can become overconfident; the uniformity tests in Section~\ref{sec:uniformity} quantify exactly this failure. Fourth, SC3R's evidence has boundaries: the pre-registered gate passes in full at $k{=}4$ on all 15 MVTec classes (seeds 0--2) and replicates at $k{=}8$ with power-gain intervals excluding zero everywhere, but the $k{=}8$ false-alarm budget is met only at $\alpha{=}0.05$---at $\alpha{=}0.10$ the pooled rate is 0.121 against a 0.120 budget with no-harm 74\%---and at both $k$ the exceedances concentrate in Gaussian noise and JPEG. SC3R also requires a multi-category deployment with matched acquisition conditions, and cross-dataset transfer of source archives (e.g., MVTec sources for VisA targets) is untested. Fifth, ECE on conformal-derived scores is strongly prevalence-sensitive, so calibration numbers on anomaly-balanced test sets must not be read as deployment reliability; we treat operational false-alarm, power, and precision metrics as primary.

The remaining extensions are scoped: SC3R with more seeds and cross-dataset source archives, randomized/smoothed few-shot p-values as an alternative route below the attainable-alpha floor, and a vision--language few-shot baseline such as WinCLIP~\cite{winclip2023} run under our audit rather than cited from reported numbers.
